%-----------------------------------------------------------------------
\subsection{Instance 5: Causal Discovery Within a Markov Equivalence Class}
\label{sec:instance-causal}
%-----------------------------------------------------------------------

We instantiate Theorem~\ref{thm:universal} for causal structure learning,
where the fundamental obstruction is Markov equivalence: distinct directed
acyclic graphs (DAGs) can be observationally indistinguishable yet disagree
on causal edge directions.

\paragraph{The explanation system.}
Define the 6-tuple $\mathcal{S}_{\mathrm{causal}} =
(\Params,\, \sH,\, Y,\, \mathsf{obs},\, \mathsf{exp},\, \incomp)$ as follows.

\begin{itemize}
  \item $\Params$ is the set of DAGs over a fixed vertex set $V$
    (the space of possible causal structures, i.e., data-generating processes).
  \item $\sH$ is the set of DAGs over $V$ used as \emph{explanations}: the
    causal structure that a discovery algorithm commits to.  The attribution
    map $\theta : \sH \to \Params$ sends each explanation DAG to the
    underlying data-generating DAG it was fitted to.
  \item $Y$ is the set of \emph{conditional independence (CI) structures}:
    each $y \in Y$ records, for every triple $(X_i, X_j \mid S)$ of
    variables, whether $X_i \perp\!\!\!\perp X_j \mid S$ holds in the
    associated distribution.  Equivalently, $Y$ indexes Markov equivalence
    classes.
  \item $\mathsf{obs} : \sH \to Y$ maps each DAG $h \in \sH$ to the
    CI structure it encodes (the set of $d$-separation statements implied by
    $h$ via the Markov condition).  Two DAGs map to the same element of $Y$
    if and only if they are \emph{Markov equivalent}.
  \item $\mathsf{exp} : \sH \to \sH$ is an \emph{orientation rule}: given
    a DAG $h \in \sH$ as fitted by an algorithm, the rule commits to a fully
    oriented causal graph.  In practice, $\mathsf{exp}$ is any causal
    discovery procedure (PC algorithm, GES, FCI, etc.) that returns a single
    DAG rather than an equivalence class.
  \item $\incomp \subseteq \sH \times \sH$ is the \emph{edge-reversal
    incompatibility}: $(h, h') \in \incomp$ if there exists a directed edge
    $X \to Y$ in $h$ and $Y \to X$ in $h'$ (or vice versa).  Two causal
    explanations are incompatible when they assert opposite orientations for
    at least one edge.
\end{itemize}

\paragraph{The Rashomon property: Markov equivalence.}
The key structural fact is that Markov equivalence classes are non-trivial
for nearly all graph structures of practical interest.

\begin{definition}[Markov Equivalence]
\label{def:markov-equiv}
Two DAGs $h, h'$ over $V$ are \emph{Markov equivalent} if they encode
the same set of conditional independence relations, i.e.,
$\mathsf{obs}(h) = \mathsf{obs}(h')$.  A Markov equivalence class
$[h]$ contains all DAGs with the same CI structure; it is represented
by a \emph{completed partially directed acyclic graph} (CPDAG).
\end{definition}

\citet{verma1991equivalence} established that two DAGs are Markov equivalent
if and only if they share the same skeleton (undirected graph) and the same
set of v-structures (immoralities $X \to Z \leftarrow Y$ with no edge between
$X$ and $Y$).  Consequently, any undirected edge in the CPDAG corresponds to
a reversible orientation: there exist two Markov-equivalent DAGs that orient
it in opposite directions.  \citet{chickering2002optimal} gave the canonical
characterization of these classes and the GES algorithm that outputs CPDAGs
rather than DAGs precisely to avoid committing to an arbitrary orientation.

\begin{proposition}[Causal Rashomon Property]
\label{prop:causal-rashomon}
$\mathcal{S}_{\mathrm{causal}}$ has the Rashomon property: for any CPDAG
containing at least one undirected edge $\{X, Y\}$, there exist DAGs
$h, h' \in \sH$ with $\mathsf{obs}(h) = \mathsf{obs}(h')$ (same CI
structure) yet $(h, h') \in \incomp$ (opposite orientations of $X$--$Y$).
\end{proposition}

\begin{proof}
Let $h$ be any DAG in the equivalence class orienting the edge as $X \to Y$,
and let $h'$ be the DAG obtained by reversing that edge to $Y \to X$ (which
exists in the equivalence class precisely because the edge is undirected in
the CPDAG).  By Markov equivalence, $\mathsf{obs}(h) = \mathsf{obs}(h')$.
By definition of $\incomp$, $(h, h') \in \incomp$.
\end{proof}

This is the formalization captured in \texttt{CausalInstance.lean} and
\texttt{CausalDiscovery.lean}: the \texttt{rashomon} field of
\texttt{causalSystem} is a concrete witness pair $(h, h')$ from
Proposition~\ref{prop:causal-rashomon} (specifically, the 3-node chain vs.\
fork).  \emph{Note:} The Lean formalization uses the broader relation
$\texttt{incompatible}(h_1, h_2) \iff h_1 \ne h_2$ rather than the
edge-reversal relation defined above.  Since $(h_1 \ne h_2)$ implies
edge-reversal incompatibility for Markov-equivalent DAGs with a reversed edge,
the impossibility under the broader relation implies the impossibility under
the narrower one \emph{a fortiori}.  The paper's edge-reversal definition is
the scientifically meaningful one; the Lean code uses the broader relation for
simplicity of formalization.

\paragraph{Instance 5: causal discovery impossibility.}

\begin{theorem}[Causal Discovery Impossibility --- Instance 5]
\label{thm:causal-impossibility}
The explanation system $\mathcal{S}_{\mathrm{causal}}$ has the Rashomon
property.  Therefore, by Theorem~\ref{thm:universal}, no orientation rule
$\mathsf{exp}$ can simultaneously be:
\begin{enumerate}
  \item \textbf{Faithful}: $\mathsf{exp}(h) = h$ for every DAG $h \in \sH$
    (the returned DAG matches the true data-generating structure);
  \item \textbf{Stable}: $\mathsf{exp}(h) = \mathsf{exp}(h')$ whenever
    $\mathsf{obs}(h) = \mathsf{obs}(h')$ (Markov-equivalent structures
    receive identical causal explanations); and
  \item \textbf{Decisive}: $\mathsf{exp}$ always commits to a single fully
    oriented DAG (never returns a partial orientation or equivalence class).
\end{enumerate}
\end{theorem}

\begin{proof}
By Proposition~\ref{prop:causal-rashomon}, $\mathcal{S}_{\mathrm{causal}}$
has the Rashomon property.  Apply Theorem~\ref{thm:universal}.
\end{proof}

{\raggedright
This is verified in Lean~4 as
\texttt{causal\_instance\_impossibility} in
\texttt{CausalInstance.lean} (three axioms), and as the
self-contained \texttt{causal\_discovery\_impossibility}
in \texttt{CausalDiscovery.lean} (zero additional axioms
beyond the \texttt{CausalOrientationProblem} structure).\par}

\paragraph{Empirical illustration.}
Three causal discovery algorithms (PC with $\alpha = 0.05$, PC with
$\alpha = 0.01$, GES) on 1{,}000 samples from the Asia network show
overall edge agreement of only 0.17 and \emph{orientation} agreement of 0.33
(Mann-Whitney $p = 4.4 \times 10^{-38}$ across 100 random seeds per condition).
At $N = 100{,}000$, orientation agreement reaches~1.0 (all algorithms agree on
every directed edge), though overall agreement is 0.625 because the three
undirected edges in the CPDAG are counted as non-agreeing by the overall
metric.  This distinction is important: the algorithms converge to the
\emph{same CPDAG}, confirming that finite-sample ambiguity (not algorithmic
limitation) drives the small-$N$ disagreement.

\paragraph{Connection to the classical Markov equivalence problem.}
Theorem~\ref{thm:causal-impossibility} is the classical Markov equivalence underdetermination
recast in the universal framework.  The three properties above correspond
exactly to the three goals that any ``complete'' causal discovery method would
need to achieve: return the true DAG (faithfulness), be invariant to
observationally equivalent parameterizations (stability), and produce a unique
answer (decisiveness).  The Markov equivalence theorem of
\citet{verma1991equivalence} has long been understood to block this
combination; our contribution is to show it is a \emph{special case} of the
single four-line proof of Theorem~\ref{thm:universal}.

\paragraph{Resolution: CPDAGs as $G$-invariant explanations.}
The uniform resolution of Section~\ref{sec:resolution} specializes to causal
discovery as follows.  Let $G = [h]$ denote the Markov equivalence class
(the symmetry group acting by edge reversals on the set of reversible edges).
The $G$-invariant explanation map aggregates over all member DAGs:
\[
  \mathsf{exp}^*(h)
  \;=\;
  \mathrm{CPDAG}([h]),
\]
the completed partially directed acyclic graph representing the equivalence
class.  This map is \textbf{faithful in expectation} (it captures all shared
causal structure), \textbf{stable} by construction (all Markov-equivalent
DAGs map to the same CPDAG), but \textbf{not decisive}: undirected edges in
the CPDAG correspond precisely to the unresolvable orientation queries.
Reporting the CPDAG is therefore the principled response to the
impossibility---trading decisiveness for faithfulness and stability, exactly
as predicted by Theorem~\ref{thm:universal}.  This justification for CPDAG
reporting is formalized in \texttt{CausalDiscovery.lean} as
\texttt{cpdag\_is\_stable}.
